\documentclass[fontsize=11pt, parskip=half]{scrartcl}
\pagestyle{plain}

\usepackage{amsmath, amssymb, amsthm, amsfonts}
\newtheorem{theorem}{Theorem}

\usepackage[dvipsnames]{xcolor}

\usepackage[draft]{commenting}
\declareauthor{leon}{LL}{blue}
\authorcommand{leon}{comment}

\usepackage[colorlinks=true, linkcolor=blue, urlcolor=blue]{hyperref}

% Solution toggle: set to true to show solutions, false to hide
\newif\ifsolutions
\solutionstrue  % comment out or set \solutionsfalse to hide

\usepackage{tcolorbox}
\tcbuselibrary{breakable}

\ifsolutions
  \newenvironment{solution}{%
    \begin{tcolorbox}[colback=black!5!white, colframe=black!50!white, 
      title={\textbf{Solution}}, fonttitle=\normalfont, breakable]
  }{\end{tcolorbox}}
\else
  \newenvironment{solution}{\expandafter\comment}{\expandafter\endcomment}
  \usepackage{verbatim} % for \comment environment
\fi

%%%%%%%%%%%%%%%%%%%%%%%%%%%



\newcommand{\mytitle}[1]{{\noindent\Large\textbf{#1}}}
\newcommand{\mysection}[1]{\textbf{\section*{#1}}}
\newcommand{\mysubsection}[2]{\romannumeral #1) #2}


%===================================
\begin{document}

\noindent\textbf{Iliad Intensive} \hfill \textbf{London Initiative for Safe AI}\\
Leon Lang  \hfill April 2026\\

\mytitle{Solomonoff Induction Exercises \hfill \today}


%===================================

\mysection{Setup and Definitions}

\leon{Motivate these exercises more: Why is the result it's building toward interesting?}

\leon{Perhaps I should name the theorem we show in this exercise sheet}

For more setup and definitions, see \href{https://www.lesswrong.com/posts/HSDumToH57nSRdLST/a-technical-introduction-to-solomonoff-induction-without-k}{the corresponding post on Solomonoff induction}.

A \textbf{semiprobability distribution} over a finite set $\mathcal{X}$ is a function $Q: \mathcal{X} \to [0, 1]$ with $Q(x) \geq 0$ and $\sum_{x \in \mathcal{X}} Q(x) \leq 1$.

A \textbf{semimeasure} over infinite binary sequences is a function $\nu: \{0, 1\}^* \to [0, 1]$ with $\nu(\epsilon) \leq 1$ and $\nu(x) \geq \nu(x0) + \nu(x1)$ for all $x \in \{0, 1\}^*$.
A semimeasure is a \textbf{measure} if equality holds for both equations.

The \textbf{universal a priori distribution} $M$ is a lower semicomputable semimeasure, which can be written in the form
\begin{equation}\label{eq:mixture_expression}
M(x_{\leq n}) = \sum_{\nu \in \mathcal{M}_{sol}} P_{ap}(\nu) \nu(x_{\leq n})
\end{equation}
for all finite strings $x_{\leq n}$, where $\mathcal{M}_{sol}$ is the set of all lower semicomputable semimeasures and $P_{ap}(\nu)$ is the \textbf{a priori probability} of $\nu$.

With $\nu_i, i = 1, 2, \dots$ an effective enumeration of all lower semicomputable discrete semimeasures and the \textbf{Solomonoff prior} $P_{sol}(i) = 2^{-K(i)}$, there also exist two constants $0 < c < C$ such that
\begin{equation}\label{eq:second:representation}
 c \cdot \sum_{i \in \mathbb{N}} P_{sol}(i) \nu_i (x_{\leq n}) < M(x_{\leq n}) < C \cdot \sum_{i \in \mathbb{N}} P_{sol}(i) \nu_i (x_{\leq n}).
\end{equation}





\mysection{Problem 1}

Let $P$ be a probability distribution over $\mathcal{X}$ and $Q$ be a semiprobability distribution over the same set. The \emph{Kullback--Leibler (KL) divergence} from $P$ to $Q$ is defined as
\[
D_{\text{KL}}\big( P \ \| \ Q \big) := \sum_{x \in \mathcal{X}} P(x) \ln \frac{P(x)}{Q(x)},
\]
where we use the conventions $0 \ln \frac{0}{q} = 0$ for any $q \geq 0$, and $p \ln \frac{p}{0} = \infty$ for $p > 0$.

Now let $(X, Y)$ be a pair of random variables taking values in $\mathcal{X} \times \mathcal{Y}$, and let $P$ and $Q$ be two joint distributions over $(X, Y)$ (with $Q$ only being a \emph{semi}probability distribution). Define the \emph{conditional KL divergence} as
\[
D_{\text{KL}}\big( P(Y \mid X) \ \| \ Q(Y \mid X) \big) := \sum_{x \in \mathcal{X}} P(x) \sum_{y \in \mathcal{Y}} P(y \mid x) \ln \frac{P(y \mid x)}{Q(y \mid x)}.
\]

Show that the KL divergence satisfies the \emph{chain rule}:
\[
D_{\text{KL}}\big( P(X, Y) \ \| \ Q(X, Y) \big) = D_{\text{KL}}\big( P(X) \ \| \ Q(X) \big) + D_{\text{KL}}\big( P(Y \mid X) \ \| \ Q(Y \mid X) \big).
\]

\begin{solution}
We expand the joint KL divergence using $P(x,y) = P(x)P(y \mid x)$ and $Q(x,y) = Q(x)Q(y \mid x)$:
\begin{align*}
D_{\text{KL}}\big( P(X, Y) \ & \| \ Q(X, Y) \big) = \sum_{x, y} P(x, y) \ln \frac{P(x, y)}{Q(x, y)} \\
&= \sum_{x, y} P(x) P(y \mid x) \ln \frac{P(x) P(y \mid x)}{Q(x) Q(y \mid x)} \\
&= \sum_{x, y} P(x) P(y \mid x) \left[ \ln \frac{P(x)}{Q(x)} + \ln \frac{P(y \mid x)}{Q(y \mid x)} \right] \\
&= \sum_{x} P(x) \ln \frac{P(x)}{Q(x)} \underbrace{\sum_{y} P(y \mid x)}_{= 1} + \sum_{x} P(x) \sum_{y} P(y \mid x) \ln \frac{P(y \mid x)}{Q(y \mid x)} \\
&= D_{\text{KL}}\big( P(X) \ \| \ Q(X) \big) + D_{\text{KL}}\big( P(Y \mid X) \ \| \ Q(Y \mid X) \big). \qedhere
\end{align*}
\end{solution}

\mysection{Problem 2}

Let $P$ be a joint probability distribution over random variables $X_1, \dots, X_n$ taking values in $\mathcal{X}_1 \times \cdots \times \mathcal{X}_n$, and let $Q$ be a \emph{semi}probability distribution over the same set of variables. Using the notation $X_{<k} := (X_1, \dots, X_{k-1})$, show by induction on $n$ (using the result of Problem~1) that
\[
D_{\text{KL}}\big( P(X_1, \dots, X_n) \ \| \ Q(X_1, \dots, X_n) \big) = \sum_{t=1}^{n} D_{\text{KL}}\big( P(X_t \mid X_{<t}) \ \| \ Q(X_t \mid X_{<t}) \big),
\]
where $P(X_1 \mid X_{<1}) := P(X_1)$ and similarly for $Q$.

\begin{solution}
We prove this by induction on $n$.

\textbf{Base case ($n = 1$):} The sum reduces to $D_{\text{KL}}\big( P(X_1) \ \| \ Q(X_1) \big)$, which is just the left-hand side.

\textbf{Induction step ($n - 1 \to n$):} Apply the chain rule from Problem~1 with $X = X_{<n}$ and $Y = X_n$:
\begin{align*}
& D_{\text{KL}}\big( P(X_1, \dots, X_n) \ \| \ Q(X_1, \dots, X_n) \big) \\
&= D_{\text{KL}}\big( P(X_{<n}) \ \| \ Q(X_{<n}) \big) + D_{\text{KL}}\big( P(X_n \mid X_{<n}) \ \| \ Q(X_n \mid X_{<n}) \big).
\end{align*}
By the induction hypothesis, the first term satisfies
\[
D_{\text{KL}}\big( P(X_{<n}) \ \| \ Q(X_{<n}) \big) = \sum_{t=1}^{n-1} D_{\text{KL}}\big( P(X_t \mid X_{<t}) \ \| \ Q(X_t \mid X_{<t}) \big).
\]
Combining both gives
\[
D_{\text{KL}}\big( P(X_1, \dots, X_n) \ \| \ Q(X_1, \dots, X_n) \big) = \sum_{t=1}^{n} D_{\text{KL}}\big( P(X_t \mid X_{<t}) \ \| \ Q(X_t \mid X_{<t}) \big). \qedhere
\]
\end{solution}


\mysection{Problem 3}

Show that for all $z > 0$,
\[
\ln(z) \geq 1 - \frac{1}{z}.
\]
\textit{Hint:} Consider the function $f(z) = \ln(z) - 1 + \frac{1}{z}$ and show that it attains its global minimum at $z = 1$.


\begin{solution}
Define $f(z) = \ln(z) - 1 + \frac{1}{z}$. We want to show $f(z) \geq 0$ for all $z > 0$. We have
\[
f'(z) = \frac{1}{z} - \frac{1}{z^2} = \frac{z - 1}{z^2}.
\]
Thus $f'(z) < 0$ for $z \in (0, 1)$ and $f'(z) > 0$ for $z > 1$, so $f$ has a global minimum at $z = 1$. Since
\[
f(1) = \ln(1) - 1 + 1 = 0,
\]
we conclude $f(z) \geq 0$ for all $z > 0$, i.e., $\ln(z) \geq 1 - \frac{1}{z}$.
\end{solution}


\mysection{Problem 4}

Let $P$ be a probability distribution over $\mathbb{B} = \{0, 1\}$ and let $Q$ be a semiprobability distribution over $\mathbb{B}$. Write $p_0 = P(0)$, $p_1 = P(1) = 1 - p_0$, $q_0 = Q(0)$, $q_1 = Q(1)$. We want to show that
\[
\sum_{x \in \mathbb{B}} \big( Q(x) - P(x) \big)^2 \leq \sum_{x \in \mathbb{B}} P(x) \ln \frac{P(x)}{Q(x)}.
\]

\mysubsection{1}{Show that the inequality holds whenever $Q(x) = 0$ for some $x \in \mathbb{B}$.}

\textit{Hint:} Distinguish the cases $P(x) > 0$ and $P(x) = 0$, and use Problem~3 for the latter.

\mysubsection{2}{Assuming $q_0, q_1 > 0$, show that it suffices to prove the inequality for the case $q_0 + q_1 = 1$, i.e., when $Q$ is a full probability distribution.}

\textit{Hint:} Define $F(q_0, q_1) := p_0 \ln \frac{p_0}{q_0} + p_1 \ln \frac{p_1}{q_1} - (q_0 - p_0)^2 - (q_1 - p_1)^2$ and show that $F$ is decreasing in $q_1$.

\mysubsection{3}{Now assume $q_0 + q_1 = 1$. Show that}
\[
2(p_0 - q_0)^2 \leq p_0 \ln \frac{p_0}{q_0} + p_1 \ln \frac{p_1}{1 - q_0}
\]
and observe this finishes the proof.

\textit{Hint:} Define $g(q_0)$ as the right-hand side minus the left-hand side, compute $g'(q_0)$, and use that $q_0(1 - q_0) \leq \frac{1}{4}$.


\begin{solution}

\textbf{Part (i):}
If $P(x) > 0$ and $Q(x) = 0$ for some $x \in \mathbb{B}$, then the right-hand side contains the term $P(x) \ln \frac{P(x)}{0} = +\infty$, so the inequality holds trivially.

If $P(x) = 0$ and $Q(x) = 0$ for some $x$, then that $x$ contributes $0$ to both sides (using the convention $0 \ln \frac{0}{0} = 0$). The other value $x' \neq x$ has $P(x') = 1$ and $Q(x') \leq 1$, so the inequality reduces to
\[
(Q(x') - 1)^2 \leq \ln \frac{1}{Q(x')}.
\]
By Problem~3 applied to $z = \frac{1}{Q(x')}$, we get $\ln \frac{1}{Q(x')} \geq 1 - Q(x')$. Since $1 - Q(x') \in [0, 1]$, we have $(1 - Q(x'))^2 \leq 1 - Q(x') \leq \ln \frac{1}{Q(x')}$.

\medskip
\textbf{Part (ii):}
Define
\[
F(q_0, q_1) := p_0 \ln \frac{p_0}{q_0} + p_1 \ln \frac{p_1}{q_1} - (q_0 - p_0)^2 - (q_1 - p_1)^2.
\]
We want to show $F(q_0, q_1) \geq 0$. We compute
\[
\frac{\partial F}{\partial q_1} = -\frac{p_1}{q_1} - 2(q_1 - p_1).
\]
This has no real roots as a function of $q_1$: setting it to zero gives $2q_1^2 - 2p_1 q_1 + p_1 = 0$, whose discriminant is $4p_1(p_1 - 2) < 0$. Since $\frac{\partial F}{\partial q_1}\big|_{q_1 = p_1} = -1 < 0$, the derivative is strictly negative for all $q_1 > 0$.

Thus $F$ is decreasing in $q_1$. For any semiprobability distribution with $q_0 + q_1 \leq 1$, increasing $q_1$ to $1 - q_0$ only decreases $F$:
\[
F(q_0, q_1) \geq F(q_0, 1 - q_0).
\]
So it suffices to show $F(q_0, 1 - q_0) \geq 0$.

\medskip
\textbf{Part (iii):}
Setting $q_1 = 1 - q_0$, we have $(q_0 - p_0)^2 + (q_1 - p_1)^2 = 2(q_0 - p_0)^2$, so we need to show
\[
g(q_0) := p_0 \ln \frac{p_0}{q_0} + p_1 \ln \frac{p_1}{1 - q_0} - 2(p_0 - q_0)^2 \geq 0
\]
for all $q_0 \in (0, 1)$. Differentiating:
\[
g'(q_0) = -\frac{p_0}{q_0} + \frac{p_1}{1 - q_0} + 4(p_0 - q_0) = (q_0 - p_0) \left[ \frac{1}{q_0(1 - q_0)} - 4 \right].
\]
Since $q_0(1 - q_0) \leq \frac{1}{4}$, the bracket is non-negative. Thus $g'(q_0) \leq 0$ for $q_0 < p_0$ and $g'(q_0) \geq 0$ for $q_0 > p_0$, so $g$ has a global minimum at $q_0 = p_0$. Since $g(p_0) = 0$, we conclude $g(q_0) \geq 0$ for all $q_0 \in (0, 1)$. \qedhere
\end{solution}






\mysection{Problem 5}

Let $\mu$ be any lower semicomputable \emph{measure} (!) on binary sequences, and assume it generates our actual universe. Let $M$ be the Solomonoff mixture distribution.

Define the cumulative expected prediction error of predicting sequences sampled by $\mu$ via $M$ as:
\[
S^{\mu}_{\infty} := \sum_{t = 1}^{\infty} \sum_{x_{<t} \in \mathbb{B}^*} \mu(x_{<t}) \sum_{x_t \in \mathbb{B}} \big( M(x_t \mid x_{<t}) - \mu(x_t \mid x_{<t}) \big)^2.
\]

Show that
\[
S_{\infty}^{\mu} \leq - \ln P_{ap}(\mu).
\]

\textit{Hint:} Apply Problem~4 to each summand to pass from squared error to KL divergence, then make the infinite sum explicit as a limit of finite sums up to $n$, then use Problem~2 to telescope the KL divergences, and finally use the mixture representation of $M$.

\begin{solution}
For each $t$ and $x_{<t}$, the conditional $\mu(\cdot \mid x_{<t})$ is a probability distribution over $\mathbb{B}$ (since $\mu$ is a measure), while $M(\cdot \mid x_{<t})$ is a semiprobability distribution over $\mathbb{B}$. Thus, by Problem~4:
\[
\sum_{x_t \in \mathbb{B}} \big( M(x_t \mid x_{<t}) - \mu(x_t \mid x_{<t}) \big)^2 \leq \sum_{x_t \in \mathbb{B}} \mu(x_t \mid x_{<t}) \ln \frac{\mu(x_t \mid x_{<t})}{M(x_t \mid x_{<t})}.
\]
Thus:
\begin{align*}
S_{\infty}^{\mu} &\leq \sum_{t=1}^{\infty} \sum_{x_{<t} \in \mathbb{B}^*} \mu(x_{<t}) \sum_{x_t \in \mathbb{B}} \mu(x_t \mid x_{<t}) \ln \frac{\mu(x_t \mid x_{<t})}{M(x_t \mid x_{<t})} \\
&= \lim_{n \to \infty} \sum_{t=1}^{n} D_{\text{KL}}\big( \mu(X_t \mid X_{<t}) \ \| \ M(X_t \mid X_{<t}) \big),
\end{align*}
where the inner sum over $x_{<t}$ weighted by $\mu(x_{<t})$ gives exactly the conditional KL divergence. By Problem~2 (the telescoped chain rule), this equals
\[
S_{\infty}^{\mu} \leq \lim_{n \to \infty} D_{\text{KL}}\big( \mu(X_1, \dots, X_n) \ \| \ M(X_1, \dots, X_n) \big) = \lim_{n \to \infty} \sum_{x_{\leq n} \in \mathbb{B}^n} \mu(x_{\leq n}) \ln \frac{\mu(x_{\leq n})}{M(x_{\leq n})}.
\]
Now, since $\mu \in \mathcal{M}_{sol}$, the mixture representation gives
\[
M(x_{\leq n}) = \sum_{\nu \in \mathcal{M}_{sol}} P_{ap}(\nu) \nu(x_{\leq n}) \geq P_{ap}(\mu) \mu(x_{\leq n}).
\]
Therefore $\frac{\mu(x_{\leq n})}{M(x_{\leq n})} \leq \frac{1}{P_{ap}(\mu)}$, and so
\[
S_{\infty}^{\mu} \leq \lim_{n \to \infty} \sum_{x_{\leq n}} \mu(x_{\leq n}) \ln \frac{1}{P_{ap}(\mu)} = -\ln P_{ap}(\mu) \cdot \lim_{n \to \infty} \underbrace{\sum_{x_{\leq n}} \mu(x_{\leq n})}_{= 1} = -\ln P_{ap}(\mu),
\]
where $\sum_{x_{\leq n}} \mu(x_{\leq n}) = 1$ since $\mu$ is a measure. \qedhere
\end{solution}


\mysection{Problem 5'}

Under the same setup as Problem~5, show that there exists a constant $\kappa \in \mathbb{R}$ (independent of $\mu$) such that
\[
S_{\infty}^{\mu} \leq K(\mu) \ln 2 + \kappa,
\]
where $K(\mu) := \min\{ K(i) : \nu_i = \mu \}$ is the Kolmogorov complexity of $\mu$ (minimized over all indices $i$ in the effective enumeration from the setup that give rise to $\mu$).

\textit{Hint:} Follow the same strategy as Problem~5, but instead of the mixture representation Equation~\eqref{eq:mixture_expression}, use the lower bound $M(x_{\leq n}) > c \cdot \sum_{i} P_{sol}(i)\, \nu_i(x_{\leq n})$ from Equation~\eqref{eq:second:representation}, together with $P_{sol}(i) = 2^{-K(i)}$.

\begin{solution}
The proof follows the same steps as Problem~5 up to the KL divergence expression:
\[
S_{\infty}^{\mu} \leq \lim_{n \to \infty} \sum_{x_{\leq n}} \mu(x_{\leq n}) \ln \frac{\mu(x_{\leq n})}{M(x_{\leq n})}.
\]
Now, by Equation~\eqref{eq:second:representation}, we have
\[
M(x_{\leq n}) > c \cdot \sum_{i \in \mathbb{N}} P_{sol}(i)\, \nu_i(x_{\leq n}) \geq c \cdot P_{sol}(i^*)\, \mu(x_{\leq n}),
\]
where $i^*$ is any index with $\nu_{i^*} = \mu$. Therefore $\frac{\mu(x_{\leq n})}{M(x_{\leq n})} < \frac{1}{c \cdot P_{sol}(i^*)}$, and so
\[
S_{\infty}^{\mu} \leq \lim_{n \to \infty} \sum_{x_{\leq n}} \mu(x_{\leq n}) \ln \frac{1}{c \cdot P_{sol}(i^*)} = -\ln c - \ln P_{sol}(i^*) = -\ln c + K(i^*) \ln 2,
\]
using $P_{sol}(i^*) = 2^{-K(i^*)}$ and $\sum_{x_{\leq n}} \mu(x_{\leq n}) = 1$. Since this holds for every $i^*$ with $\nu_{i^*} = \mu$, we may take the minimum over all such indices to obtain
\[
S_{\infty}^{\mu} \leq K(\mu) \ln 2 + \kappa,
\]
where $\kappa = -\ln c$. \qedhere
\end{solution}

\leon{Throughout, there are some subtle trivial compatibilities between semimeasures and semiprobability distributions that I use, e.g. that semimeasures are semiprobability distributions on fixed-length sequences, and that they give rise to CONDITIONALS that are actual probability distributions. Maybe I could mention that or make these their own exercises? Hhm.}

\mysection{Interpretation}


Problems~5 and 5' show that the \emph{cumulative} expected prediction error $S_\infty^\mu$ is finite. Since $S_\infty^\mu$ is a sum of non-negative terms, this immediately implies that the per-step expected prediction error
\[
d_t := \sum_{x_{<t} \in \mathbb{B}^*} \mu(x_{<t}) \sum_{x_t \in \mathbb{B}} \big( M(x_t \mid x_{<t}) - \mu(x_t \mid x_{<t}) \big)^2
\]
converges to zero as $t \to \infty$. In other words, $M$'s predictions become indistinguishable from $\mu$'s predictions on average over histories.



\end{document}